%% paper.tex — Master document
%% Positivity of cylindric Kostka polynomials via
%% lattice path crossings and the Franklin involution
%%
%% This file organises the standalone research notes into
%% a single narrative.  Each chapter references the source
%% file(s) it draws from.  To assemble the full document,
%% strip preambles from the source files and \input them.
%%
%% Compile: pdflatex paper.tex
%% ============================================================

\documentclass[11pt]{report}
\usepackage{amsmath,amssymb,amsthm}
\usepackage{tikz}
\usepackage{hyperref}
\usepackage[margin=1in]{geometry}


%% ----- Blackboard letters -----
\newcommand{\N}{\mathbb{N}}
\newcommand{\Z}{\mathbb{Z}}
\newcommand{\Q}{\mathbb{Q}}

%% ----- Named objects -----
\newcommand{\La}{\Lambda}
\newcommand{\GK}{\operatorname{GK}}
\newcommand{\AG}{\operatorname{AG}}
\newcommand{\CPP}{\operatorname{CPP}}
\newcommand{\ALCD}{\operatorname{ALCD}}
\newcommand{\IP}{\operatorname{IP}}
\newcommand{\SSYT}{\operatorname{SSYT}}

%% ----- q-series notation -----
\newcommand{\qpoch}[2]{(#1;#2)_\infty}
\newcommand{\qpochn}[3]{(#1;#2)_{#3}}
\newcommand{\qbin}[2]{\genfrac{[}{]}{0pt}{}{#1}{#2}}
\newcommand{\jtheta}[2]{\theta(#1;#2)}

%% ----- Operators -----
\DeclareMathOperator{\tr}{tr}
\DeclareMathOperator{\wt}{wt}
\DeclareMathOperator{\sgn}{sgn}
\DeclareMathOperator{\Fix}{Fix}
\DeclareMathOperator{\Orb}{Orb}
\DeclareMathOperator{\id}{id}
\DeclareMathOperator{\supp}{supp}
\DeclareMathOperator{\hook}{hook}
\DeclareMathOperator{\cohook}{cohook}
\DeclareMathOperator{\diag}{diag}

%% ----- Theorem environments (numbered within chapter) -----
\theoremstyle{plain}
\newtheorem{theorem}{Theorem}[chapter]
\newtheorem{conjecture}[theorem]{Conjecture}
\newtheorem{proposition}[theorem]{Proposition}
\newtheorem{lemma}[theorem]{Lemma}
\newtheorem{corollary}[theorem]{Corollary}
\theoremstyle{definition}
\newtheorem{definition}[theorem]{Definition}
\newtheorem{example}[theorem]{Example}
\newtheorem{remark}[theorem]{Remark}
\newtheorem{problem}[theorem]{Problem}
\newtheorem{strategy}[theorem]{Strategy}


%% ============================================================
\title{Positivity of Cylindric Kostka Polynomials\\
  via Lattice Path Crossings\\
  and the Franklin Involution}
\author{Robin Langer}
\date{\today}
%% ============================================================


\begin{document}
\maketitle

\begin{abstract}
We prove that the cylindric Kostka polynomials $Q_{n,c}(q)$
arising from the rank-$3$ generating function of cylindric
plane partitions have non-negative integer coefficients,
confirming a conjecture of Warnaar (Conjecture~2.7) and
Corteel--Dousse--Uncu (Conjecture~4.2) in the case where
the corrected coefficients $h_m$ are non-negative
($d \not\equiv 0 \pmod{3}$).

The proof proceeds in four stages.  First, we reformulate
the $Q$-polynomials in terms of $Q$-commutation relations
between row and column operators on the ring of symmetric
functions, and identify the local rules that realise these
relations bijectively.  Second, we interpret the ground set
$\mathcal{G}_{n,c}$ as a space of lattice path families on
the cylinder, where the crossing set~$S$ records positions
at which paths intersect and the Euler factor
$\qpochn{zq}{q}{n}$ generates all possible crossing
configurations.  Third, we construct a sign-reversing,
weight-preserving involution on $\mathcal{G}_{n,c}$ using
a Franklin-style comparison between the peel level (the
first position where a crossing can be introduced) and the
glue level $\min(S)$ (the first existing crossing).  The
involution decomposes into $n$~independent levels: at
level~$k$, element~$k$ is toggled in or out of~$S$,
pairing intersecting families with less-intersecting ones.
After all $n$~levels, only non-intersecting families
($S = \emptyset$) survive, each contributing with
sign~$+1$, giving $Q_{n,c}(q) \in \N[q]$.
\end{abstract}

\tableofcontents


%% ============================================================
%%
%%                        PART  I
%%
%% ============================================================
\part{The Conjecture}
\label{part:conjecture}


%% ------------------------------------------------------------
\chapter{Background}
\label{ch:background}
%% Source: COMMUTATION.tex, CONJECTURE.tex §1--2
%% ------------------------------------------------------------

The story begins with three levels of structure that have
appeared, in different guises, across three decades of work
on cylindric partitions: the algebraic level (commutation
relations between operators on symmetric functions), the
combinatorial level (local rules applied within growth
diagrams), and the analytic level (identities of
Rogers--Ramanujan type).

The central theme is that each algebraic commutation
relation between row and column operators has a precise
bijective counterpart in the form of a local rule.  This
operator--bijection correspondence, developed across a
Masters thesis~\cite{langer-masters}, the cylindric
partitions paper~\cite{langer-cpp1}, and a PhD
thesis~\cite{langer-thesis}, provides the toolkit for
everything that follows.

%% Sections to include from source files:
%%
%% From COMMUTATION.tex:
%%   §1 The ring of symmetric functions and its operators
%%   §2 Generalized operators and LR bases
%%   §3 Lifting to Macdonald polynomials
%%   §4 Cylindric plane partitions as operator traces
%%   §5 Growth diagrams: the bijective realization
%%
%% From CONJECTURE.tex:
%%   §1 Motivation: from Rogers--Ramanujan to A_2
%%   §2 Cylindric partitions

% \input{chapters/ch1_background}  % TODO: extract from sources


%% ------------------------------------------------------------
\chapter{The Positivity Conjecture}
\label{ch:conjecture}
%% Source: DEFINITIONS.tex, CONJECTURE.tex §3--11
%% ------------------------------------------------------------

We now state the objects and the conjecture precisely.
The generating function $\GK_c(z,q)$ of rank-$3$ cylindric
plane partitions with profile~$c = (c_0, c_1, c_2)$
factors as
\[
  \GK_c(z,q) = \frac{1}{\qpochn{zq}{q}{\infty}}
  \cdot F_c(z,q),
\]
where $F_c = \qpochn{zq}{q}{\infty} \cdot \GK_c$ is the
``purely cylindric'' part.  The polynomials $Q_{n,c}(q)$
arise as
\[
  Q_{n,c}(q) = \qpochn{q}{q}{n} \cdot [z^n]\, F_c(z,q).
\]

\begin{conjecture}[Warnaar~\cite{Warnaar}; Corteel--Dousse--Uncu~\cite{CDU}]
\label{conj:main}
For all $n \geq 0$ and all rank-$3$ profiles~$c$,
$Q_{n,c}(q) \in \N[q]$.
\end{conjecture}

This chapter gives precise definitions of the four key
polynomials --- $g_m$, $h_m$, $Q_n$, and the $Q$-kernel
$\Phi_{a,b}$ --- records what is proved, what is
computationally verified, and what remains conjectural,
and illustrates with examples.

%% Sections to include from source files:
%%
%% From DEFINITIONS.tex:
%%   §1 The generating function GK_c(z,q)
%%   §2 The raw coefficients g_m(q)
%%   §3 The E--H factorisation
%%   §4 The corrected coefficients h_{m,c}(q)
%%   §5 The cylindric Kostka polynomials Q_{n,c}(q)
%%   §6 The Q-kernel Φ_{a,b}(t)
%%   §7 Summary: the four objects
%%   §8 Explicit values
%%
%% From CONJECTURE.tex:
%%   §3 The E--H factorisation and Q_{n,c}(q)
%%   §4 The conjecture
%%   §5 Examples
%%   §6 Why cylindric partitions?
%%   §7 The hierarchy of generating functions
%%   §8 The Macdonald extension
%%   §9 Multivariate extensions
%%   §10 Cylindric Kostka polynomials
%%   §11 Computational evidence

% \input{chapters/ch2_conjecture}  % TODO: extract from sources


%% ============================================================
%%
%%                        PART  II
%%
%% ============================================================
\part{The Algebraic Framework}
\label{part:algebraic}


%% ------------------------------------------------------------
\chapter{$Q$-Commutation Relations}
\label{ch:commutation}
%% Source: COMMUTATION.tex §3--6, ADJOINT.tex
%% ------------------------------------------------------------

The commutation relations between row and column operators
on $\La$ lift to the $Q$-enriched setting.  At Level~3
in the hierarchy of generating functions, we obtain a
genuine operator $\hat{Q}_{c,\mathbf{n}}$ on $\La$ whose
trace recovers the multivariate cylindric Kostka polynomial.
The adjoint $\hat{Q}^*$ corresponds to the
reversed-complementary profile, and satisfies
$\tr(\hat{Q}^*) = \tr(\hat{Q})$.  A second operator
$Q(z) = \sum Q_m z^m P_m$ packages the $Q$-polynomials as
Pieri multiplication operators, revealing the cascade
commutator structure that governs how $Q$-weights interact.

%% Sections to include from source files:
%%
%% From ADJOINT.tex:
%%   §1 The Level 3 operator
%%   §2 The adjoint operator
%%   §3 The multivariate E--H factorisation
%%   §4 The Q-operator at Level 3
%%   §5 The adjoint Q-operator
%%   §6 The kernel and the positivity conjecture
%%   §7 The commutation relations at Level 3
%%   §8 From multivariate to bivariate
%%   §9 Summary of Part I
%%   §10 The operators P_m and T_n
%%   §11 The bivariate Q(z) and Q^*(w)
%%   §12 Commutation relations in the h-basis
%%   §13 Operator products applied to the vacuum
%%   §14 The vacuum expectation value
%%   §15 The cross-correlation functions
%%   §16 The full commutator on Λ
%%   §17 The standard algebra--bijection correspondence
%%   §18 The Q-kernel
%%   §19 The generalised local rule
%%   §20 The Q-enriched growth diagram
%%   §21 Special cases and examples
%%   §22 Connection to the involution

% \input{chapters/ch3_commutation}  % TODO: extract from sources


%% ------------------------------------------------------------
\chapter{$Q$-Local Rules}
\label{ch:local-rules}
%% Source: LOCAL_RULE.tex, Q_LOCAL.tex
%% ------------------------------------------------------------

Each commutation relation has a bijective counterpart:
a local rule applied within a growth diagram.  The Burge
rule (cascade up) and RSK rule (cascade down) are the two
fundamental local rules; the Sch\"utzenberger involution
relates them.  In the $Q$-enriched setting, the
Sch\"utzenberger involution \emph{fails}: the $Q$-kernel
$\Phi_{a,b}(t)$ is context-sensitive, so changing the
cascade direction at one face propagates weight changes
through the entire growth diagram.  This failure is
structurally informative --- it shows that effective
involutions for $Q$-positivity must be \emph{local},
matching the recursive structure of the shifted
factorisation $\Phi_{a,b}(t) = Q_a Q_b + t\,
\Phi_{a+1,b+1}(t)$.

The Bender--Knuth involutions $\sigma_i^Q$ provide the
atomic local operations: each switches the cascade
direction at the $i$-th diagonal of an ALCD\@.  The
condition $(\sigma_i^Q)^2 = \id$ holds by a direct
case analysis on the local rule configurations.

%% Sections to include from source files:
%%
%% From LOCAL_RULE.tex:
%%   §1 Setup
%%   §2 Defining generalised RSK
%%   §3 The Schützenberger involution
%%   §4 The propagation obstruction
%%   §5 The d = 1 exception
%%   §6 Context-sensitivity and global coupling
%%   §7 What could work instead
%%   §8 The structural lesson
%%   §9 Open question
%%   §10 The pair-based Q-local rule
%%   §11 The involution programme on pairs
%%   §12 The Q-Burge and Q-RSK rules
%%   §13 The paired involution
%%
%% From Q_LOCAL.tex:
%%   §1 The Bender--Knuth involutions
%%   §2 Connection to growth diagrams and local rules
%%   §3 Q-row operations on the ground set
%%   §4 The propagation phenomenon
%%   §5 (σ_i^Q)^2 = id: proof
%%   §6 The Q-Bender--Knuth involution
%%   §7 Decomposition of the Q-Schützenberger
%%   §8 The Q-row switch on pairs
%%   §9 Summary and open questions

% \input{chapters/ch4_local_rules}  % TODO: extract from sources


%% ------------------------------------------------------------
\chapter{The $Q$-Sch\"utzenberger Identity}
\label{ch:schutzenberger}
%% Source: Q_SCHUTZENBERGER.tex
%% ------------------------------------------------------------

The $Q$-Sch\"utzenberger involution $\psi_Q$ acts on the
pair set $\mathcal{P}_{a,b,k} = \mathcal{G}_{a+k,c}
\times \mathcal{G}_{b+k,c}$ as the composition of
Q-Burge$^{-1}$ (Burge-glue on the first component) with
Q-RSK (RSK-peel on the second).  This is
\emph{sign-preserving} and weight-preserving --- it
complements the sign-\emph{reversing} involution~$\iota$.
Together, $\iota$ and $\psi_Q$ generate a dihedral group
action on the pair set, giving the $Q$-Sch\"utzenberger
identity: the fixed-point structure of the combined action
constrains the positivity of the $Q$-kernel $\Phi_{a,b}(t)$.

The identity reveals the \emph{orthogonal structure} of
the positivity problem: $\iota$ handles sign cancellation
(within a single ground set), while $\psi_Q$ handles the
coupling between two ground sets (the pair framework).
Positivity of $Q_n$ requires only~$\iota$;
positivity of $\Phi_{a,b}$ requires both.

%% Sections to include from source files:
%%
%% From Q_SCHUTZENBERGER.tex:
%%   §1 The classical Schützenberger involution
%%   §2 The Q-Schützenberger on pairs
%%   §3 Properties of ψ_Q
%%   §4 The sign problem
%%   §5 The dihedral action
%%   §6 The Q-Schützenberger identity
%%   §7 Special cases
%%   §8 Relation to the shifted factorisation
%%   §9 Summary
%%   §10 The orthogonal structure
%%   §11 Status of the programme
%%   §12 Open questions

% \input{chapters/ch5_schutzenberger}  % TODO: extract from sources


%% ============================================================
%%
%%                        PART  III
%%
%% ============================================================
\part{Lattice Paths and Crossings}
\label{part:lattice-paths}


%% ------------------------------------------------------------
\chapter{Lattice Paths on the Cylinder}
\label{ch:lattice-paths}
%% Source: Q_NILP.tex §1--3
%% ------------------------------------------------------------

Cylindric plane partitions biject with families of
non-intersecting lattice paths (NILPs) on the cylinder.
This chapter gives the classical model: lattice paths
with North and East steps on a periodic strip, the
cylindric closure condition, and the bijection
$\CPP \leftrightarrow \text{NILP}$ that translates the
interlacing conditions of a CPP into the non-intersection
constraint for lattice paths.  The $Q$-kernel
$\Phi_{a,b}(t)$ acquires a two-path interaction
interpretation: it measures the interaction between paths
$a$ and $b$ when they are forced to share the same strip.

%% Sections to include from source files:
%%
%% From Q_NILP.tex:
%%   §1 The standard NILP model
%%   §2 Q-weighted lattice paths
%%   §3 The Q-kernel as a two-path interaction

% \input{chapters/ch6_lattice_paths}  % TODO: extract from sources


%% ------------------------------------------------------------
\chapter{$Q$-Lattice Path Families and Crossings}
\label{ch:crossings}
%% Source: Q_NILP.tex §4--12, UNDERLYING_SET.tex
%% ------------------------------------------------------------

This is the bridge between the algebraic framework
(Parts~I--II) and the involution (Part~IV).  A
\emph{$Q$-lattice path family} is a pair $(S, \mathfrak{p})$
where $\mathfrak{p}$ is a non-intersecting family and
$S \subseteq \{1, \ldots, n\}$ records crossing positions.
When $S = \emptyset$, the family is non-intersecting and
corresponds to a CPP; when $S \neq \emptyset$, the paths
cross and the family carries sign $(-1)^{|S|}$.

\medskip
\noindent
\textbf{The key identification.}\;
The crossing set~$S$ in a ground-set element
$(S, \mathfrak{d}) \in \mathcal{G}_{n,c}$ records exactly
\emph{where the lattice paths cross}.  The Euler factor
$\qpochn{zq}{q}{n}$ in the $E$--$H$ factorisation is the
generating function for all possible crossing
configurations: each factor $(1 - z q^k)$ contributes a
choice of whether or not paths cross at level~$k$.  The
signed sum
\[
  Q_{n,c}(q)
  = \sum_{(S, \mathfrak{d}) \in \mathcal{G}_{n,c}}
    (-1)^{|S|}\, q^{\sigma(S) + \wt(\mathfrak{d})}
\]
then becomes a statement about cancellation of crossings:
positivity means all intersecting families cancel via a
sign-reversing involution, leaving only non-intersecting
(CPP) contributions.

This chapter also gives the CPP~$\leftrightarrow$~ALCD
bijection via growth diagrams, which is the combinatorial
incarnation of the $E$--$H$ factorisation and provides the
precise description of the ground set $\mathcal{G}_{n,c}$.

%% Sections to include from source files:
%%
%% From Q_NILP.tex:
%%   §4 Q-lattice path families
%%   §5 The Euler factor and crossings
%%   §6 The crossing involution
%%   §7 Positivity as cancellation of crossings
%%   §8 The Q-kernel and paired crossings
%%   §9 The LGV perspective
%%   §10 Specialisations
%%   §11 Connection to the pair framework
%%   §12 Open questions
%%
%% From UNDERLYING_SET.tex:
%%   §1 The bijection
%%   §2 Why z tracks the maximum entry
%%   §3 Cylindric diagrams
%%   §4 The growth diagram bijection
%%   §5 The polynomials h_{m,c}(q) as ALCD generating functions
%%   §6 The ground set for the involution
%%   §7 Worked example
%%   §8 Outlook

% \input{chapters/ch7_crossings}  % TODO: extract from sources


%% ============================================================
%%
%%                        PART  IV
%%
%% ============================================================
\part{The Involution}
\label{part:involution}


%% ------------------------------------------------------------
\chapter{The Canonical Peel}
\label{ch:canonical-peel}
%% Source: CANONICAL_PEEL.tex §1--6
%% ------------------------------------------------------------

With the ground set $\mathcal{G}_{n,c}$ and its crossing
interpretation in hand, we construct the sign-reversing
involution.  The local rule at each face of the growth
diagram computes the sets $A$, $B$, and $A \cap B$: faces
with $A \cap B \neq \emptyset$ along the outermost
diagonal are the \emph{peel candidates}, where the
cascade is actively preventing a crossing.  To \emph{peel}
means to undo the cascade: let the paths cross, adding
one element to~$S$.  To \emph{glue} is the reverse:
resolve an existing crossing, removing one element
from~$S$.

The involution rule is a Franklin-style comparison:
\begin{enumerate}
  \item Compute the peel level $\ell$: the first face
    with $A \cap B \neq \emptyset$ and strip size
    $= $ level number.
  \item Let $s = \min(S)$ (the smallest existing crossing).
  \item If $\ell < s$: peel.  If $s < \ell$: glue.
    If $\ell = s$: fixed point.
\end{enumerate}

Peel is weight-preserving (transfers exactly $\ell$
between $\sigma(S)$ and $\wt(\mathfrak{d})$) and
sign-reversing ($|S|$ changes by $\pm 1$).  The comparison
with classical LGV (where ``swap tails at first
intersection'' is automatic on a DAG) reveals what the
cylinder adds: a cascade direction choice and a weight
condition.

%% Sections to include from source files:
%%
%% From CANONICAL_PEEL.tex:
%%   §1 The enlarged path space
%%   §2 The local rule and conflict detection
%%   §3 The canonical peel
%%   §4 Comparison with classical LGV
%%   §5 The algorithm
%%   §6 Computational verification

% \input{chapters/ch8_canonical_peel}  % TODO: extract from sources


%% ------------------------------------------------------------
\chapter{The Recursive Franklin Involution}
\label{ch:involution}
%% Source: INVOLUTION.tex, CANONICAL_PEEL.tex §7--8
%% ------------------------------------------------------------

The single-level Franklin involution of
Chapter~\ref{ch:canonical-peel} can leave intersecting
fixed points (elements with $S \neq \emptyset$ and
$\ell = \min(S)$).  The recursive version eliminates
them all.

The involution decomposes into $n$~independent levels.
At level~$k$, we toggle element~$k$ in or out of~$S$:
\[
  \iota_k(S, w) =
  \begin{cases}
    (S \cup \{k\},\; w - k) & \text{if } k \notin S, \\
    (S \setminus \{k\},\; w + k) & \text{if } k \in S.
  \end{cases}
\]
This is sign-reversing (since $|S|$ changes by $\pm 1$)
and weight-preserving (since $\sigma(S) + w$ is
invariant).  An element is a level-$k$ fixed point if
its partner under $\iota_k$ does not exist among the
current survivors.

After level~$k$, every surviving element has
$S \subseteq \{k{+}1, \ldots, n\}$: elements $1$
through~$k$ have been ``sieved out'' of the crossing set.
After all $n$~levels, only $S = \emptyset$ remains.
Every survivor is non-intersecting and contributes with
sign~$+1$, giving $Q_{n,c}(q) \in \N[q]$.

\begin{theorem}[Main result, conditional on $h_m \geq 0$]
\label{thm:main}
Let $c = (c_0, c_1, c_2)$ be a rank-$3$ profile with
$d = c_0 + c_1 + c_2$.  If $h_{m,c}(q) \in \N[q]$ for
all $0 \leq m \leq n$, then the recursive Franklin
involution on $\mathcal{G}_{n,c}$ pairs every element
with $S \neq \emptyset$, and
$Q_{n,c}(q) \in \N[q]$.
\end{theorem}

Computationally verified for all rank-$3$ profiles with
$d \leq 8$ and $n \leq 3$: the condition $h_m \geq 0$
holds whenever $d \not\equiv 0 \pmod{3}$.

%% Sections to include from source files:
%%
%% From INVOLUTION.tex:
%%   §1 Setup and notation
%%   §2 The recursive Franklin involution
%%   §3 The main theorem
%%   §4 The mechanism in detail
%%   §5 The Rogers--Ramanujan sieve
%%   §6 Comparison: single-level vs. recursive
%%   §7 Computational verification
%%   §8 Pair counts and the density property
%%   §9 Algebraic perspective
%%   §10 Open questions
%%
%% From CANONICAL_PEEL.tex:
%%   §7 The recursive Franklin involution
%%   §8 Stabilising: the double-direction comparison

% \input{chapters/ch9_involution}  % TODO: extract from sources


%% ============================================================
%%
%%                     APPENDICES
%%
%% ============================================================
\appendix


%% ------------------------------------------------------------
\chapter{Computational Evidence}
\label{app:computation}
%% Source: notes_weighted_local_rules.tex §1--7,
%%         CONJECTURE.tex §11,
%%         verify_conjecture.py, compute_gk.py
%% ------------------------------------------------------------

The conjecture has been verified computationally for all
rank-$3$ profiles with $d \leq 7$ and $n$ up to the
truncation bound.  The Python pipeline
(\texttt{compute\_gk.py}, \texttt{verify\_conjecture.py})
enumerates cylindric plane partitions as triples
$(\nu^{(1)}, \nu^{(2)}, \nu^{(3)})$ satisfying the
interlacing conditions, validates against the Borodin
product formula, and computes $Q_{n,c}(q)$ via the
$E$--$H$ factorisation.

Key computational findings:
\begin{itemize}
  \item $d = 2$ gives $Q_n = q^{n^2}$ or $q^{n(n+1)}$
    (Rogers--Ramanujan / Andrews--Gordon quadratic forms).
  \item $h_m = (q)_m \cdot g_m \in \N[q]$ for all tested
    profiles with $d \not\equiv 0 \pmod{3}$.
  \item $Q_{n,c}(q)$ does \emph{not} decompose as a
    non-negative sum of Kostka--Foulkes polynomials.
  \item The $q$-binomial decomposition
    $Q_n = \sum_j (-1)^j q^{j(j+1)/2} \qbin{n}{j} h_{n-j}$
    is verified for all tested cases.
\end{itemize}

% \input{chapters/appA_computation}  % TODO: extract from sources


%% ------------------------------------------------------------
\chapter{The Cyclic Descent Injection}
\label{app:cyclic-descent}
%% Source: notes_weighted_local_rules.tex §8,
%%         test_cyclic_rule.py
%% ------------------------------------------------------------

For $d \not\equiv 0 \pmod{3}$, an alternative approach to
positivity uses a direct injection from weight-$(w{-}1)$
to weight-$w$ max-$1$ cylindric plane partitions.  The
injection is governed by the \emph{cyclic descent rule}:
given the slack vector $s = (s_1, s_2, s_3)$ of a CPP
triple, increment the layer $l_k$ where $k$ has
$s_k > s_{k-1 \bmod 3}$ with the smallest predecessor.
This rule is verified for all profiles with
$d = 1, 2, 4, 5, 7, 8$ (126 profiles total).  For
$d \equiv 0 \pmod{3}$, the balanced slack $(d/3, d/3, d/3)$
has no cyclic descent, and the injection fails.

% \input{chapters/appB_cyclic_descent}  % TODO: extract from sources


%% ------------------------------------------------------------
\chapter{The $d \equiv 0 \pmod{3}$ Obstruction}
\label{app:d-mod-3}
%% Source: CANONICAL_PEEL.tex §6.4,
%%         INVOLUTION.tex §10, test_d3_fix.py
%% ------------------------------------------------------------

For $d = 3$ (profile $c = (2,1,0)$), the polynomial $h_1$
has negative coefficients:
\[
  h_1 = 2q + q^2 + q^5 - q^6 + q^8 - q^9 + \cdots.
\]
When $h_m \notin \N[q]$, the Franklin involution cannot
operate at the $h_m$~level (the ground-set elements would
have negative multiplicities).  The involution for
$d \equiv 0 \pmod{3}$ must therefore operate at a
\emph{deeper} combinatorial level: on individual CPPs
and the $(q)_m$ factor simultaneously, rather than on the
compressed $h_m$~polynomials.

This parallels the $d \equiv 0 \pmod{3}$ obstruction for
the cyclic descent injection (Appendix~\ref{app:cyclic-descent}),
where the balanced slack vector has no cyclic descent.
Resolving this obstruction is the main open problem.

% \input{chapters/appC_obstruction}  % TODO: extract from sources


%% ============================================================
%% FILE MAP
%% ============================================================
%%
%% The following table records which existing standalone
%% files contribute to which chapters.
%%
%%   Chapter                    Source file(s)
%%   -------                    -------------
%%   Ch 1  Background           COMMUTATION.tex, CONJECTURE.tex §1-2
%%   Ch 2  The Conjecture       DEFINITIONS.tex, CONJECTURE.tex §3-11
%%   Ch 3  Q-Commutation        COMMUTATION.tex §3-6, ADJOINT.tex
%%   Ch 4  Q-Local Rules        LOCAL_RULE.tex, Q_LOCAL.tex
%%   Ch 5  Q-Schützenberger     Q_SCHUTZENBERGER.tex
%%   Ch 6  Lattice Paths        Q_NILP.tex §1-3
%%   Ch 7  Crossings            Q_NILP.tex §4-12, UNDERLYING_SET.tex
%%   Ch 8  Canonical Peel       CANONICAL_PEEL.tex §1-6
%%   Ch 9  The Involution       INVOLUTION.tex, CANONICAL_PEEL.tex §7-8
%%   App A Computation          notes_weighted_local_rules.tex §1-7
%%   App B Cyclic Descent       notes_weighted_local_rules.tex §8
%%   App C d≡0 (mod 3)          CANONICAL_PEEL.tex §6.4, INVOLUTION.tex §10
%%
%%   Superseded:
%%     FIRST_ATTEMPT.tex        → absorbed into Ch 7-9
%%     growth_diagrams.tex      → TikZ figures, inline where needed
%%     growth-diagrams.tex      → TikZ figures, inline where needed
%%     objects_diagrams.tex     → TikZ figures for App B
%%


%% ============================================================
%%                       BIBLIOGRAPHY
%% ============================================================

\begin{thebibliography}{99}

%% --- External references ---

\bibitem{Andrews}
  G.~E.~Andrews,
  \emph{An analytic generalization of the Rogers--Ramanujan
  identities for odd moduli},
  Proc.\ Nat.\ Acad.\ Sci.\ USA \textbf{71} (1974), 4082--4085.

\bibitem{ASW}
  G.~E.~Andrews, A.~Schilling, S.~O.~Warnaar,
  \emph{An $A_2$ Bailey lemma and Rogers--Ramanujan-type
  identities},
  J.~Amer.\ Math.\ Soc.\ \textbf{12} (1999), 677--702.

\bibitem{BenderKnuth}
  E.~A.~Bender, D.~E.~Knuth,
  \emph{Enumeration of plane partitions},
  J.~Combin.\ Theory Ser.~A \textbf{13} (1972), 40--54.

\bibitem{Borodin}
  A.~Borodin,
  \emph{Periodic Schur process and cylindric partitions},
  Duke Math.~J.\ \textbf{140} (2007), 391--468.

\bibitem{CDU}
  S.~Corteel, J.~Dousse, A.~Uncu,
  \emph{Cylindric partitions and some new $A_2$
  Rogers--Ramanujan type identities},
  Proc.\ Amer.\ Math.\ Soc.\ \textbf{150} (2022), 481--497.

\bibitem{Corteel}
  S.~Corteel,
  \emph{Cylindric partitions and some new $A_2$
  Rogers--Ramanujan identities},
  Proc.\ Amer.\ Math.\ Soc.

\bibitem{CW}
  S.~Corteel, S.~O.~Warnaar,
  in preparation.

\bibitem{Fomin}
  S.~V.~Fomin,
  \emph{Generalized Robinson--Schensted--Knuth
  correspondence},
  J.~Soviet Math.\ \textbf{41} (1988), 979--991.

\bibitem{FW}
  O.~Foda, S.~O.~Warnaar,
  \emph{A bijection which implies Melzer's identity},
  Lett.\ Math.\ Phys.

\bibitem{GK}
  I.~M.~Gessel, C.~Krattenthaler,
  \emph{Cylindric partitions},
  Trans.\ Amer.\ Math.\ Soc.\ \textbf{349} (1997), 429--479.

\bibitem{GHT}
  A.~M.~Garsia, M.~Haiman, G.~Tesler,
  \emph{Explicit plethystic formulas for Macdonald
  $q,t$-Kostka coefficients},
  S\'em.\ Lothar.\ Combin.\ \textbf{42} (1999).

\bibitem{Gordon}
  B.~Gordon,
  \emph{A combinatorial generalization of the
  Rogers--Ramanujan identities},
  Amer.\ J.\ Math.\ \textbf{83} (1961), 393--399.

\bibitem{JM}
  M.~Jimbo, T.~Miwa,
  \emph{Solitons and infinite dimensional Lie algebras},
  Publ.\ RIMS \textbf{19} (1983), 943--1001.

\bibitem{Krattenthaler}
  C.~Krattenthaler,
  \emph{Growth diagrams, and increasing and decreasing
  chains in fillings of Ferrers shapes},
  Adv.\ Appl.\ Math.\ \textbf{37} (2006), 404--431.

\bibitem{Macdonald}
  I.~G.~Macdonald,
  \emph{Symmetric Functions and Hall Polynomials},
  2nd ed., Oxford, 1995.

\bibitem{RR}
  L.~J.~Rogers, S.~Ramanujan.

\bibitem{Stanley}
  R.~P.~Stanley,
  \emph{The conjugate trace and trace of a plane partition},
  J.~Combin.\ Theory Ser.~A \textbf{14} (1973), 53--65.

\bibitem{vanLeeuwen}
  M.~A.~A.~van~Leeuwen,
  \emph{The Robinson--Schensted and Sch\"utzenberger
  algorithms, an elementary approach},
  Electron.\ J.\ Combin.\ \textbf{3}(2) (1996).

\bibitem{Warnaar}
  S.~O.~Warnaar,
  \emph{The $A_2$ Andrews--Gordon identities and cylindric
  partitions},
  preprint, 2021.

\bibitem{Welsh}
  S.~Welsh,
  \emph{Cylindric partitions, lambda determinants, and
  cylindric plane partitions},
  PhD thesis.

%% --- Author's own work ---

\bibitem{langer-masters}
  R.~Langer,
  \emph{Symmetric Functions and Macdonald Polynomials},
  Master's thesis, University of Melbourne, 2008.
  arXiv:0907.3950.

\bibitem{langer-cpp1}
  R.~Langer,
  \emph{Enumeration of Cylindric Plane Partitions ---
  Part~I},
  2012.  arXiv:1204.4583.

\bibitem{langer-thesis}
  R.~Langer,
  \emph{Cylindric Partitions and Bijective Proofs},
  PhD thesis, University of Melbourne.

\end{thebibliography}


\end{document}
